% Table 5.9 (matheuristic vs exact) and the parameter paragraph of 5.8, 2026-09-27.
The three parameters of Section~\ref{subsec:ils} are set once for every instance and every experiment that follows: the cap divisor $D = 3$, so that a shake removes at most $c = \lceil k/3 \rceil$ targets, the length $L_r = 20$ of the restricted candidate list of the reordering operators, and the shake limit $S = 100$. The look-ahead of Equation~\eqref{eq:shake-value} scores the next six removals of the enumeration. Table~\ref{tab:theta} reports the choice of $D$.

The shake limit is read from the traces rather than tuned by repeated runs, since every run records the shake count at each improvement of the best found solution. On the fourteen runs of Table~\ref{tab:matheuristic-vs-exact} the longest gap between two successive improvements is $88$ shakes, on PR11~(48), and at most $46$ elsewhere, so $S = 100$ is the smallest round value that cuts no improvement; $S = 50$ would have lost $0.95\%$ on PR11 and saved a third of the shakes, and $S = 25$ up to $6.1\%$ on PR15~(240). The restricted list is a trade, as Section~\ref{subsec:proposal} states. Against the unrestricted reordering sets, $L_r = 20$ raises the objective by $7.78\%$ on PR15~(240), $4.62\%$ on PR10~(288) and $0.99\%$ on C104~(100), where it multiplies the shakes a run affords from $116$, $89$ and $36$ to $691$, $311$ and $381$; on PR11~(48), whose $33$ scheduled targets leave a reordering set that is affordable in full, it costs $0.70\%$, of which the unrestricted set recovers $0.56$ points at $313$ shakes in $43$\,s against $146$ shakes in $17$\,s. On the instances whose windows fix the order the list is not binding.

\caption{Greedy (R4) vs.\ single-start matheuristic vs.\ MISOCP at $\eta = 1$, the design of Section~\ref{sec:heuristic}. The matheuristic is not capped by the clock and ends after $S = 100$ consecutive shakes without an improvement (Section~\ref{subsec:ils}), with $D = 3$ and $L_r = 20$. $t_{\mathrm{best}}$ is the wall-clock time at which the ILS reaches its best solution and Run is the time the whole run takes, the MISOCP gap is from Table~\ref{tab:scalability} (bold marks gap $\le 0.5\%$), and $\Delta = (f^{\star} - f_{\mathrm{best}}^{\mathrm{ILS}})/f^{\star}$ is negative when the ILS beats the MISOCP incumbent.}\label{tab:matheuristic-vs-exact}
\begin{tabular}{l r rrr r r c r}
\toprule
& \multicolumn{1}{c}{Greedy (R4)} & \multicolumn{3}{c}{Single-start ILS ($S = 100$)} & \multicolumn{3}{c}{Exact MISOCP} & \\
\cmidrule(lr){2-2}\cmidrule(lr){3-5}\cmidrule(lr){6-8}
Instance & $f(R_4)$ & $f_{\mathrm{best}}^{\mathrm{ILS}}$ & $t_{\mathrm{best}}$ (s) & Run (s) & $f^{\star}$ & time (s) & gap (\%) & $\Delta$ (\%) \\
\midrule
R101 (50)       & 4\,643.08 & 11\,902.02 & 0.1 & 1.9 & \textbf{11\,921.13} & \textbf{7.7} & \textbf{0.00} & $+0.16$ \\
R101 (100)      & 8\,830.33 & 22\,884.44 & 0.2 & 4.2 & \textbf{22\,886.91} & \textbf{40.9} & \textbf{0.00} & $+0.01$ \\
R1\_2\_1 (200)  & 5\,758.63 & 16\,586.35 & 9.2 & 17.8 & \textbf{16\,954.02} & \textbf{99.6} & \textbf{0.00} & $+2.17$ \\
C101 (50)       & 2\,432.97 & 8\,591.92 & 1.5 & 33.4 & 8\,594.11 & 3\,600.0 & 2.61 & $+0.03$ \\
C101 (100)      & 3\,776.53 & 11\,324.41 & 110.2 & 147.3 & 11\,328.36 & 3\,600.1 & 32.19 & $+0.03$ \\
C1\_2\_1 (200)  & 2\,583.30 & 11\,146.32 & 78.0 & 138.7 & 10\,747.43 & 3\,600.4 & 87.46 & $\mathbf{-3.71}$ \\
RC1\_2\_1 (200) & 4\,198.18 & 17\,973.77 & 70.1 & 90.3 & 16\,746.34 & 3\,600.2 & 209.93 & $\mathbf{-7.33}$ \\
R102 (100)      & 13\,854.06 & 30\,443.78 & 42.9 & 59.2 & 27\,941.61 & 3\,600.1 & 108.51 & $\mathbf{-8.95}$ \\
R104 (100)      & 15\,524.42 & 38\,105.79 & 39.9 & 57.7 & 25\,377.52 & 3\,600.1 & 271.42 & $\mathbf{-50.16}$ \\
C104 (100)      & 9\,670.98 & 18\,414.60 & 217.4 & 304.3 & 16\,468.21 & 3\,600.2 & 32.01 & $\mathbf{-11.82}$ \\
RC104 (100)     & 19\,444.49 & 36\,231.33 & 17.2 & 27.8 & 27\,439.42 & 3\,600.1 & 289.22 & $\mathbf{-32.04}$ \\
PR11 (48)       & 1\,858.09 & 6\,473.95 & 6.9 & 17.3 & 5\,401.01 & 3\,600.0 & 78.53 & $\mathbf{-19.87}$ \\
PR15 (240)      & 2\,471.30 & 19\,639.68 & 609.4 & 704.8 & 6\,005.65 & 3\,600.8 & 713.79 & $\mathbf{-227.02}$ \\
PR10 (288)      & 3\,113.48 & 16\,696.23 & 128.4 & 196.6 & 5\,377.05 & 3\,600.9 & 984.58 & $\mathbf{-210.51}$ \\
\bottomrule
\end{tabular}
\end{table}
